% ============================================================
% 第 15 章 其他常用感应电机（教材第四篇，印刷页 167--174）
% 节名：第一节 单相感应电动机；第二节 直线感应电机；第三节 感应发电机
% ============================================================
\chapter{其他常用感应电机}

三相感应电动机之外，还有单相感应电动机（家用电器、小功率场合）、直线感应电机（直线运动场合）和感应发电机（风力发电等）。它们的基本原理仍由第 10 章"脉振磁动势 = 两个反向旋转磁动势"这一条推出。

\section{单相感应电动机}

\subsection{单相绕组的磁场}
单相感应电动机定子上只有一相工作绕组（单相），通入单相交流电流后产生的是\textbf{脉振磁动势}，其轴线固定、幅值随时间脉振，而不是旋转磁场。由第 10 章第二节，
\begin{equation}
f_{\phi1}(x,t)=\frac{F_{\phi1}}{2}\sin\left(\omega t-\frac{\pi}{\tau}x\right)
+\frac{F_{\phi1}}{2}\sin\left(\omega t+\frac{\pi}{\tau}x\right),
\end{equation}
即一个脉振磁场可分解为\textbf{两个幅值相等（各为 $F_{\phi1}/2$）、转速相同（均为 $n_1$）、转向相反}的旋转磁场。

\subsection{为什么没有起动转矩}
设转子转速为 $n$，则转子对正转磁场的转差率为
\begin{equation}
s_+=\frac{n_1-n}{n_1}=s,
\end{equation}
对反转磁场的转差率为
\begin{equation}
s_-=\frac{n_1+n}{n_1}=2-s .
\end{equation}
正转磁场产生正向转矩 $T_+$，反转磁场产生反向转矩 $T_-$。在起动瞬间 $n=0$，$s_+=s_-=1$，两磁场对称，故 $T_+=T_-$，合成转矩
\begin{equation}
T=T_++T_-=0 .
\end{equation}
\textbf{单相感应电动机没有起动转矩}，不能自行起动——这是它区别于三相感应电动机的根本特点。

若用外力使转子转起来（$n>0$），则 $s_+<1$、$s_->1$，正转转矩大于反转转矩，合成转矩为正，电机可继续加速到接近同步转速运行；转向取决于初始转动方向，故单相电动机"往哪边拨就往哪边转"。

\subsection{起动方法}
解决起动问题的思路只有一个：\textbf{人为制造一个在空间上不同相位的第二套绕组电流，使气隙磁场"转"起来}。

\subsubsection{分相式（电阻分相）}
在定子上另加一套\textbf{起动绕组}，其在空间上与工作绕组相差 $90^\circ$ 电角度；起动绕组用细导线（电阻大、电抗小），使其电流与工作绕组电流在时间上接近 $90^\circ$ 相位差，从而产生椭圆形旋转磁场、产生起动转矩。起动后由离心开关切除起动绕组。

\subsubsection{电容分相式}
在起动绕组中串入电容器，使起动绕组电流\textbf{超前}工作绕组电流接近 $90^\circ$，可获得接近圆形的旋转磁场，起动转矩较大、起动电流较小。若电容长期接入，则称\textbf{电容运转式}，可提高运行功率因数与效率。

\subsubsection{罩极式}
在凸极磁极的一部分套上短路铜环（罩极环），环中感应电流使被罩部分的磁通滞后，形成"移行磁场"而产生起动转矩。结构最简单、最可靠，但起动转矩小、效率低，多用于小风扇。

\begin{center}
\begin{tabular}{@{}p{2.2cm}p{3.6cm}p{3.2cm}p{3.4cm}@{}}
\toprule
起动方式 & 原理 & 起动转矩 & 特点 \\
\midrule
电阻分相 & 起动绕组电阻大，电流相位差 $<90^\circ$ & 中等 & 简单便宜，离心开关切除 \\
电容分相 & 串电容，相位差接近 $90^\circ$ & 较大 & 可做成电容运转式提高功率因数 \\
罩极式 & 短路环使磁通移相 & 小 & 结构最简单，效率低 \\
\bottomrule
\end{tabular}
\end{center}

\corepoint{单相感应电动机是"脉振磁动势"最典型的应用：正、反两个旋转磁场叠加 $\Rightarrow$ 起动时两转矩相等 $\Rightarrow$ 无起动转矩；转动起来后 $s_+\ne s_-$ $\Rightarrow$ 有净转矩。记住 $s_+=s$、$s_-=2-s$ 这一对式子，几乎所有单相电机的定性题都能推出答案。}

\section{直线感应电机}

直线感应电机的定子被"剖开铺平"，成为\textbf{初级}（原边）；转子相应成为\textbf{次级}（副边），可以是金属板（如铝板、钢板）。初级通三相电流产生\textbf{行波磁场}（而非旋转磁场），在次级中感应电流，产生直线方向的电磁推力，使次级（或初级）做直线运动。

主要特点：
\begin{itemize}
  \item 直接产生直线运动，不需旋转—直线转换机构（如丝杠、齿条），结构简单、无磨损、可靠性高；
  \item 存在\textbf{端部效应}（初级两端磁场畸变）和\textbf{边缘效应}（横向磁场不均），使效率与功率因数低于旋转电机；
  \item 气隙较大（机械上必须留出间隙），励磁电流大、功率因数低。
\end{itemize}
应用于地铁与磁浮列车牵引、传送带、电磁弹射、电梯等。

\section{感应发电机}

感应发电机由原动机拖动到\textbf{超过同步转速}（$s<0$）运行，定子接电网（或接电容器自励），向电网输出有功功率。特点：
\begin{itemize}
  \item \textbf{需要无功励磁}：感应发电机没有独立的励磁绕组，靠电网（或并联电容器）提供励磁电流，故输出有功的同时吸收无功，功率因数滞后；
  \item \textbf{不能独立建压}：除非并联电容器自励（自励感应发电机），否则必须并网运行；
  \item \textbf{结构简单、免维护、无需同步装置}：这是它在风力发电（尤其是早期定桨距风机）中被大量采用的原因；
  \item 转速变化范围小（$s$ 绝对值很小），效率略低于同步发电机。
\end{itemize}

双馈感应发电机（绕线型转子经变频器接电网）可在变速范围内运行并独立调节有功、无功，是现代变速风电机组的主流机型。
